\scanpage{4}
i.e. un hom de schémas $\operatorname{Spec}K' - \operatorname{Spec}K$, on a un morphisme de Topos correspondant
\[
(1)\qquad B_{K'}\unsure{-|\longrightarrow}B_{K^{\deleted{\unsure{\prime}}}}
\]
associé à un hom de groupoïdes fond.
\[
(2)\qquad \Pi_{K'}\unsure{\longrightarrow|-}\Pi_{K^{\deleted{\unsure{\prime}}}}.
\]
Ceci s'explicite en disant qu'un \deleted{clôture} objet de $\Pi_{K'}$ (i.e. pt. de $B_{K'}$, ou \unsure{rev.} universel de $B_{K'}$, \deleted{)} ou clôture \deleted{sép.}\inserted{$\bar K'$} de $K'$) en définit un de $\Pi_{K^{\deleted{\unsure{\prime}}}}$ (associé, en prenant $\bar K^{\deleted{\unsure{\prime}}}=$ clôture \inserted{alg.}\deleted{sép} de $K$ dans $\bar K'$) et pour deux pts correspondants, on a un hom des groupes fondamentaux correspondants, qui s'interprète p.ex. comme
\[
(3)\qquad E_{\bar K'/K'}\unsure{-|\longrightarrow}E_{\bar K^{\deleted{\unsure{\prime}}}/K^{\deleted{\unsure{\prime}}}}
\]
et qui permet de reconstituer l'hom. de topos comme une ``restriction des scalaires''.

\margin{\unsure{1}}
L'image de (3) est le sous-groupe fermé de $E_{\bar K/K}$ qui correspond à la sous-extension
\margin{3}
